%! Sinusoidal Function Transformations — Unit 6, Lesson 6.7 — Group Activity
\documentclass[10pt]{article}
\usepackage{precalculus-article}
\usepackage{precalculus-boxes}
\newcommand{\TallMath}[1]{$\displaystyle #1\rule[-1.4em]{0pt}{3.2em}$}

\begin{document}

\pageheader{Unit 6, Lesson 6.7}{Group Activity}
\namepartnerperiod

\vspace{0.10in}

%----------------------------------------------------------------------
% Context
%----------------------------------------------------------------------
\begin{scenariobox}[Context: Transformations of Sinusoidal Functions]{sky}
The general form $f(\theta)=a\sin(b(\theta+c))+d$ encodes four independent
features of a sinusoidal function. In this activity your group will decode
these parameters from equations, connect them to transformations, and reason
about how changes to each parameter affect the function. Choose \textbf{one
tier} based on your readiness level. Everyone in your group should be prepared
to explain your reasoning.
\end{scenariobox}

\vspace{0.08in}

%----------------------------------------------------------------------
% Tier R
%----------------------------------------------------------------------
\begin{tcolorbox}[colback=white, colframe=black!40,
    title=\textbf{Tier R --- Reinforce},
    fonttitle=\bfseries, arc=2mm, left=3mm, right=3mm, top=2mm, bottom=2mm]

\textbf{Directions:} For each sinusoidal function below (no phase shift),
complete the parameter table. Then answer the questions.

\vspace{4pt}
\renewcommand{\arraystretch}{1.6}
\begin{tabular}{|l|c|c|c|c|c|}
\hline
\textbf{Function} & \textbf{$a$} & \textbf{$b$} & \textbf{Amplitude} & \textbf{Period} & \textbf{Midline} \\
\hline
$y = 3\sin(\theta) + 1$ & & & & & \\
\hline
$y = -\sin(2\theta)$ & & & & & \\
\hline
$y = \tfrac{1}{2}\sin(\pi\theta) - 3$ & & & & & \\
\hline
$y = 5\cos(\theta) + 2$ & & & & & \\
\hline
$y = -4\cos(3\theta) - 1$ & & & & & \\
\hline
\end{tabular}

\vspace{8pt}
\textbf{1.}\quad For $y = 3\sin(\theta)+1$, what are the maximum and minimum output values?

\vspace{4pt}
Maximum $= $ \underline{\hspace{2cm}} \qquad Minimum $= $ \underline{\hspace{2cm}}

\vspace{6pt}
\textbf{2.}\quad For $y = -\sin(2\theta)$, how does the negative sign affect the graph
compared to $y = \sin(2\theta)$?

\writelines{2}

\vspace{6pt}
\textbf{3.}\quad Use the midline and amplitude of $y = \tfrac{1}{2}\sin(\pi\theta) - 3$
to find the maximum and minimum values. Show your reasoning.

\vspace{4pt}
Maximum $= $ midline $+$ amplitude $= $ \underline{\hspace{1cm}} $+$ \underline{\hspace{1cm}} $=$ \underline{\hspace{1.5cm}}

\vspace{4pt}
Minimum $= $ midline $-$ amplitude $= $ \underline{\hspace{1cm}} $-$ \underline{\hspace{1cm}} $=$ \underline{\hspace{1.5cm}}

\vspace{6pt}
\textbf{4.}\quad Write a sinusoidal function with amplitude 6, period $\pi$,
midline $y = -2$, and no phase shift.

\vspace{4pt}
$f(\theta) = $ \underline{\hspace{6cm}}

\end{tcolorbox}

\vspace{0.08in}

%----------------------------------------------------------------------
% Tier A
%----------------------------------------------------------------------
\begin{tcolorbox}[colback=white, colframe=black!40,
    title=\textbf{Tier A --- Apply},
    fonttitle=\bfseries, arc=2mm, left=3mm, right=3mm, top=2mm, bottom=2mm]

\textbf{Directions:} For each function, factor the argument (if needed),
then identify all four parameters. Match each function to one of the four
graph cards provided by your teacher (record the letter).

\vspace{6pt}
\textbf{Function 1:} $f(\theta) = 2\sin\!\left(3\!\left(\theta + \dfrac{\pi}{6}\right)\right) - 1$

\vspace{4pt}
Amplitude: \underline{\hspace{1.5cm}} \quad Period: \underline{\hspace{1.5cm}} \quad
Phase shift: \underline{\hspace{2.5cm}} \quad Midline: \underline{\hspace{1.5cm}} \quad Graph card: \underline{\hspace{0.8cm}}

\vspace{8pt}
\textbf{Function 2:} $g(\theta) = -3\cos(2\theta + \pi) + 4$

\vspace{4pt}
First, factor the argument: $2\theta + \pi = 2\!\left(\theta + \underline{\hspace{1.2cm}}\right)$

\vspace{4pt}
Amplitude: \underline{\hspace{1.5cm}} \quad Period: \underline{\hspace{1.5cm}} \quad
Phase shift: \underline{\hspace{2.5cm}} \quad Midline: \underline{\hspace{1.5cm}} \quad Graph card: \underline{\hspace{0.8cm}}

\vspace{8pt}
\textbf{Function 3:} $h(\theta) = 4\sin\!\left(\dfrac{1}{2}\theta - \dfrac{\pi}{4}\right)$

\vspace{4pt}
First, factor the argument: $\dfrac{1}{2}\theta - \dfrac{\pi}{4} = \dfrac{1}{2}\!\left(\theta - \underline{\hspace{1.2cm}}\right)$

\vspace{4pt}
Amplitude: \underline{\hspace{1.5cm}} \quad Period: \underline{\hspace{1.5cm}} \quad
Phase shift: \underline{\hspace{2.5cm}} \quad Midline: \underline{\hspace{1.5cm}} \quad Graph card: \underline{\hspace{0.8cm}}

\vspace{8pt}
\textbf{Function 4:} $k(\theta) = \cos\!\left(2\theta - \dfrac{\pi}{3}\right) + 2$

\vspace{4pt}
Factor: \underline{\hspace{4cm}}

\vspace{4pt}
Amplitude: \underline{\hspace{1.5cm}} \quad Period: \underline{\hspace{1.5cm}} \quad
Phase shift: \underline{\hspace{2.5cm}} \quad Midline: \underline{\hspace{1.5cm}} \quad Graph card: \underline{\hspace{0.8cm}}

\vspace{8pt}
\textbf{Discussion:} Functions 2 and 4 are cosine functions. Could each be rewritten as
a sine function? Explain what additional phase shift would be needed.

\writelines{2}

\end{tcolorbox}

\vspace{0.08in}

%----------------------------------------------------------------------
% Tier E
%----------------------------------------------------------------------
\begin{tcolorbox}[colback=white, colframe=black!40,
    title=\textbf{Tier E --- Extend},
    fonttitle=\bfseries, arc=2mm, left=3mm, right=3mm, top=2mm, bottom=2mm]

\textbf{Directions:} Factor each argument, identify all parameters, rewrite
as a sine function, and write a complete verbal transformation description in
correct algebraic order.

\vspace{6pt}
\textbf{Function A:} $p(\theta) = -3\cos(2\theta - \pi/3) + 1$

\vspace{4pt}
\textbf{Step 1.} Factor: $2\theta - \pi/3 = 2\!\left(\theta - \underline{\hspace{1.5cm}}\right)$

\vspace{4pt}
\textbf{Step 2.} Parameters: amplitude \underline{\hspace{1.5cm}}, period \underline{\hspace{1.5cm}},
phase shift \underline{\hspace{2cm}}, midline \underline{\hspace{1.5cm}}, reflected \underline{\hspace{1cm}}.

\vspace{4pt}
\textbf{Step 3.} Rewrite as a sine function using $\cos\theta = \sin(\theta + \pi/2)$:
\vspace{2pt}
$p(\theta) = -3\sin\!\Bigl(2\Bigl(\theta - \underline{\hspace{1cm}} + \underline{\hspace{1cm}}\Bigr)\Bigr) + 1
= -3\sin\!\Bigl(2\Bigl(\theta - \underline{\hspace{1.5cm}}\Bigr)\Bigr) + 1$

\vspace{4pt}
\textbf{Step 4.} Verbal description (in correct transformation order --- horizontal
scaling/shifting before vertical scaling/shifting):

\writelines{3}

\vspace{8pt}
\textbf{Function B:} $q(\theta) = 5\sin(3\theta + \pi/2) - 2$

Complete Steps 1--4 for $q(\theta)$.

\vspace{4pt}
Step 1 (factor): \underline{\hspace{6cm}}

\vspace{4pt}
Step 2 (parameters): amplitude \underline{\hspace{1.5cm}}, period \underline{\hspace{1.5cm}},
phase shift \underline{\hspace{2cm}}, midline \underline{\hspace{1.5cm}}.

\vspace{4pt}
Step 3 (sine form --- already sine, no change needed). Notice: $\sin(\theta + \pi/2) = \cos\theta$.
Write $q$ as a cosine function:
$q(\theta) = $ \underline{\hspace{5cm}}

\vspace{4pt}
Step 4 (verbal description): \writelines{3}

\vspace{6pt}
\textbf{Challenge:} Two sinusoidal functions $f$ and $g$ have the same amplitude and
period. Must they be the same function? Explain what other information is needed to
uniquely determine a sinusoidal function and how many parameters are required.

\writelines{3}

\end{tcolorbox}

\end{document}
